%&pdftex
%%%%%%%%%%%%%%%%%%%%%%%%%%%%%%%%%%%%%%%%%%%%%%%%%%%%%%%%%%%%%%%%%%%%%%%%%%%%%%%
% ULTIMA MODIFICACION: 18 DE FEBRERO DE 1992. PARA EL CONGRESO PO's DE GRANADA.
% MAS ADELANTE, ENVIADO A J. Comput. Appl. Math.
%%%%%%%%%%%%%%%%%%%%%%%%%%%%%%%%%%%%%%%%%%%%%%%%%%%%%%%%%%%%%%%%%%%%%%%%%%%%%%%

\magnification = 1200
\font \peq = cmr8
\font \peqbf = cmbx8

\font\msymnormal=msbm10
\font\msympeq=msbm7
\font\msymmuypeq=msbm5
\newfam\numeros
\textfont\numeros=\msymnormal
\scriptfont\numeros=\msympeq
\scriptscriptfont\numeros=\msymmuypeq
\def\num{\fam\numeros\msymnormal}
\def\N{{\num N}}
\def\R{{\num R}}
\font\tipo = cmr12 at 10pt
\tipo

\def \proof {{\it Proof.} }
\def \qed {\vrule height 5 pt width 5 pt depth 0 pt}

\centerline{\bf ENDPOINT WEAK BOUNDEDNESS OF SOME}
\centerline{\bf POLYNOMIAL EXPANSIONS}

\bigskip

\centerline{Jos\'e J. Guadalupe
   \footnote{$^1$}{\peq Dpto. de Matem\'aticas. Universidad de Zaragoza. 50009
Zaragoza. Spain.}, Mario P\'erez $^1$, Francisco J. Ruiz $^1$ and 
Juan L. Varona
   \footnote{$^2$}{\peq Dpto. de Matem\'atica Aplicada. Colegio Universitario de
La Rioja. 26001 Logro\~no. Spain.}}
\footnote{}{\peq 1991 Mathematics Subject Classification: 42C10.}
\footnote{}{\peq Research supported by DGICYT under grant PB89-0181-C02-02}
\footnote{}{\peq AN UPDATED VERSION OF THIS MANUSCRIPT WERE PUBLISHED IN:
J.\ Comput.\ Appl.\ Math.\ {\peqbf 49} (1993), 93--102.}

\bigskip

{\narrower\narrower\noindent{\bf Abstract. }{\sl Let $w(x) =
(1-x)^\alpha (1+x)^\beta$ on $[-1,1]$, $\alpha, \beta \geq -1/2$
and for each function $f$ let $S_nf$ be the $n$-th expansion in the
corresponding orthonormal polynomials. We show that the operators $f
\longrightarrow u S_n(u^{-1} f)$ are not of weak $(p,p)$-type,
where $u$ is another Jacobi weight and $p$ is an endpoint of the
interval of mean convergence. The same result is shown for
expansions associated to measures of the form $d\nu = w(x)dx +
\sum_{i=1}^k M_i \delta_{a_i}$.}\smallskip}

\bigskip

{\noindent \bf \S 1. Introduction and main results.}

\medskip

Let $\mu$ be a positive measure on $\R$ with infinitely many points of increase
and such that all the moments
$$
   \int_{\R} x^nd\mu \qquad (n = 0,1,\dots)
$$
exist. Let $\{P_n\}_{n \geq 0}$ stand for the corresponding orthonormal
polynomials. For $f \in L^1(d\mu)$, let $S_nf$ denote the $n$-th partial sum of 
the
orthonormal Fourier expansion of $f$ in $\{P_n\}_{n \geq 0}$:
$$
   S_n(f,x) = \int_{\R} f(y) K_n(x,y) d\mu(y), \qquad
   K_n(x,y) = \sum_{k=1}^n P_k(x) P_k(y).
$$

The problem of the uniform boundedness of the partial sum operators $S_n$ in
weighted $L^p$ spaces, that is,
$$
   \Vert u S_nf \Vert _{L^p(d\mu)} \leq C \Vert u f \Vert _{L^p(d\mu)}
   \quad \forall n \geq 0, \ \forall f \in L^p(u^pd\mu) \eqno(1)
$$
has been completely solved only in some specific cases (this boundedness
implies, in rather general situations, the $L^p$ convergence of $S_nf$ to $f$).
For example, Badkov gave in [3] necessary and sufficient conditions for $(1)$
when $d\mu$ and $u$ are generalized Jacobi weights (earlier results can be
found in [22], [24], [21], [15]). Orthogonal Hermite and Laguerre series were
studied by Askey and Wainger ([1]) and Muckenhoupt ([16], [17]).

Let us consider the case of a Jacobi weight on the interval $[-1,1]$, that is,
$d\mu = w(x) dx$,
$$
   w(x)=(1-x)^{\alpha}(1+x)^{\beta}
$$
and let $1 < p < \infty$. If $\alpha, \beta \geq -1/2$, then (see [15])
$$
   \Vert S_nf \Vert _{L^p(w)} \leq C \Vert f \Vert _{L^p(w)}
   \quad \forall n \geq 0, \ \forall f \in L^p(w) \eqno(2)
$$
if and only if $p$ belongs to the open interval $(p_0, p_1)$, where
$$
   p_0 = {4(\alpha+1) \over 2\alpha+3}, \qquad
   p_1 = {4(\alpha+1) \over 2\alpha+1}
$$
when $\alpha \geq \beta$ (and the analogous formulas with $\alpha$ replaced
by $\beta$ if $\beta \geq \alpha$).

If both $\alpha, \beta > -1/2$, the authors proved (see [6]) that the $n$-th
partial sum operators are not of weak $(p,p)$-type when $p$ is an endpoint of
the interval of mean convegence. In theorem 1 we extend this result to the
weighted case $f \longrightarrow u S_n(u^{-1} f)$, where $u$ is also a Jacobi
weight, $u(x) = (1-x)^a (1+x)^b, \ a, b \in \R$. Now, the weighted uniform
boundedness (1) holds (see [15]) if and only if
$$
   \eqalign{
   \vert a+(\alpha+1)({1 \over p} - {1 \over 2}) \vert
   < \min \{{1 \over 4}, {\alpha+1 \over 2}\},
   \cr
   \vert b+(\beta+1)({1 \over p} - {1 \over 2}) \vert
   < \min \{{1 \over 4}, {\beta+1 \over 2}\}.
   \cr} \eqno (3)
$$

Let us state our first result.

\medskip

\proclaim{Theorem 1}. Let $\alpha, \beta \geq -1/2$,
$w(x)=(1-x)^\alpha(1+x)^\beta$, $u(x)=(1-x)^a(1+x)^b$, $1<p<\infty$. Let $S_n$
be the partial sum operators associated to $w$. If there exists a constant
$C>0$ such that for every $f \in L^p(u^pw)$ and for every $n \geq 0$
$$
   \Vert uS_nf \Vert _{L^p_*(w)} \leq	 C \Vert uf \Vert _{L^p(w)},
$$
then the inequalities
$$
   \vert a+(\alpha+1)({1 \over p}-{1 \over 2}) \vert < {1 \over 4}, \qquad
   \vert b+(\beta+1) ({1 \over p}-{1 \over 2}) \vert < {1 \over 4}
$$
are verified.

On the other hand, we also study the weak boundedness of the operators $S_n$
associated to a measure $d\nu = d\mu + \sum_{i=1}^k M_i \delta_{a_i}$,
where $\mu\{a_i\} = 0$. In the particular case of a Jacobi weight and two mass
points on $1$ and $-1$, the corresponding orthonormal polynomials were studied
by Koornwinder in [10] from the point of view of differential equations (see
also [4], [2], [11], [12]). The authors have found (see [7]) some estimates for
the orthonormal polynomials and kernels relative to this type of measures.

In this context, let us consider the polynomial expansion associated to a
measure $d\nu = w(x)dx + \sum_{i=1}^k M_i \delta_{a_i}$, where
$w(x)=(1-x)^\alpha(1+x)^\beta$, $M_i>0$ and take $u(x)=(1-x)^a(1+x)^b$ for $x
\not= a_i$, $0<u(a_i)<\infty$. With this notation, we can state the following
result.

\medskip

\proclaim{Theorem 2}. Let $\alpha, \beta \geq -1/2$, $1<p<\infty$. Then, there
exists a constant $C>0$ such that
$$
   \Vert uS_nf \Vert _{L^p_*(d\nu)} \leq C \Vert uf \Vert _{L^p(d\nu)}
   \quad \forall f \in L^p(u^pd\nu), \ \forall n \geq 0,
$$
if and only if the inequalities
$$
   \vert a+(\alpha+1)({1 \over p}-{1 \over 2}) \vert < {1 \over 4}, \qquad
   \vert b+(\beta+1)({1 \over p}-{1 \over 2}) \vert < {1 \over 4}
$$
are verified.

\bigskip

{\noindent \bf \S 2. Preliminary lemmas.}

\medskip

\noindent A basic tool in the study of Fourier series on the interval $[-1,1]$
is Pollard's decomposition of the kernels $K_n(x,t)$ (see [22], [15]): if
$\{P_n\}_{n\geq 0}$ is the sequence of polynomials orthonormal with respect to
$w(x)dx$ and $\{Q_n\}_{n\geq 0}$ is the sequence of polynomials relating to
\hbox{$(1-x^2)w(x)dx$}, then
$$
   K_n(x,t) = r_nT_{1,n}(x,t) + s_nT_{2,n}(x,t) + s_nT_{3,n}(x,t),
$$
where
$$
   T_{1,n}(x,t) = P_{n+1}(x) P_{n+1}(t),
$$
$$
   T_{2,n}(x,t) = (1-t^2) {P_{n+1}(x) Q_n(t) \over x-t},
$$
$$
   T_{3,n}(x,t) = (1-x^2) {P_{n+1}(t) Q_n(x) \over t-x}
$$
and $\{r_n\}$, $\{s_n\}$ are bounded sequences. In fact, for any measure $\mu$
on $[-1,1]$ with $\mu'>0$ $a.e.$ (in particular, for $w(x)dx$),
$$
   \vbox{\halign{\hfil #  &  \qquad \hfil # \hfil \cr
   $\displaystyle\lim_{n\to\infty} r_n = -1/2$,
   & $\displaystyle\lim_{n\to\infty} s_n = 1/2$ \cr
   }}
$$
(this can be deduced from [22] and [23] or [13]). Therefore, we can write
$$
   S_nf = r_n W_{1,n}f + s_n W_{2,n}f - s_n W_{3,n}f,
$$
where
$$
   W_{1,n}f(x) = P_{n+1}(x) \int_{-1}^1 P_{n+1}(t) f(t) w(t) dt,
$$
$$
   W_{2,n}f(x) = P_{n+1}(x) H((1-t^2) Q_n(t) f(t) w(t), x)
$$
and
$$
   W_{3,n}f(x) = (1-x^2) Q_n(x) H(P_{n+1}(t) f(t) w(t), x),
$$
$H$ being the Hilbert transform on the interval $[-1,1]$. Thus, the study of
$S_n$ can be reduced to that of $W_{i,n}$ $(i=1,2,3)$.

The boundedness of the Hilbert transform can be stated in terms of Muckenhoupt's
$A_p$ classes of weights (see [9] and [19]; throughout this paper, the Hilbert
transform, as well as the $A_p$ classes, are taken on the interval $[-1,1]$): if
$u$ is a weight on $[-1,1]$ and $1<p<\infty$, then $u \in A_p$ if and only if
$H$ is a bounded operator in $L^p(u)$, with a constant which depends only on the
$A_p$ constant of $u$.

Concerning mixed weak norm inequalities for the Hilbert transform, we can state
the following property, which can be proved in the same way as in theorem 3 of
[18]: assume that $u_1(x), u_2(x), v(x) \geq 0$, $1<p<\infty$ and there is a
constant $C>0$ such that
$$
   \Vert u_2 Hg \Vert_{L^p_*(u_1)} \leq C \Vert g \Vert_{L^p(v)}
   \quad \forall g \in L^p(v);
$$
then, there exists another constant $B>0$ which depends only on $C$, such that
for every interval $I$
$$
   \Vert u_2 \chi_I \Vert_{L^p_*(u_1)} \left(
   \int_{-1}^1 {v(x)^{-1/(p-1)} \over (\vert I \vert + \vert x-x_I \vert)^q} dx
   \right)^{1/q} \leq B, \eqno (4)
$$
$x_I$ being the centre of $I$ and $1/p + 1/q = 1$.

The polynomials $P_n$ satisfy the estimate
$$
   \vert P_n(x) \vert \leq C (1-x)^{-(2\alpha+1)/4}
   (1+x)^{-(2\beta +1)/4} \quad \forall n, \forall x \in [-1,1] \eqno (5)
$$
with a constant $C>0$ independent of $x$ and $n$. A similar estimate is
verified by $Q_n$, with $\alpha+1$ and $\beta+1$ instead of $\alpha$ and
$\beta$:
$$
   \vert Q_n(x) \vert \leq C (1-x)^{-(2\alpha+3)/4}
   (1+x)^{-(2\beta +3)/4} \quad \forall n, \forall x \in [-1,1]. \eqno (6)
$$
Thus, the following easy result will be useful.

\medskip

\proclaim{Lemma 3}. Let $r \in \R$. Then, $\vert x \vert^r \in A_p([-1,1])
\Longleftrightarrow -1 < r < p-1$.

\medskip

The same property holds if we replace $x$ by $x-a$, with $a \in [-1,1]$. Even
more, it is not difficult to show that in order to see whether a finite product
of this type of expressions belongs to $A_p$, we only need to check the
above inequalities for each factor separately.

\medskip

We will eventually need to show that some of the operators are not of
strong or weak type. In this sense, this lemma (see [14]) will be used:

\proclaim{Lemma 4}. Let ${\rm supp}\ d\alpha = [-1,1]$, $\alpha'>0$ a.e. in
$[-1,1]$, and $0<p \leq \infty$. There exists a constant $C>0$ such that if $g$
is a Lebesgue-measurable function on $[-1,1]$, then
$$
   \Vert \alpha'(x)^{-1/2}(1-x^2)^{-1/4} \Vert _{L^p(\vert g \vert^pdx)} \leq
   C \liminf_{n \to \infty} \Vert P_n \Vert _{L^p(\vert g \vert^pdx)}.
$$

\medskip

There is a weak version of this property: it is a consequence of Kolmogorov's
condition (see [5], lemma V.2.8, p. 485) and the previous lemma.

\proclaim{Lemma 5}. Let ${\rm supp}\ d\alpha = [-1,1]$, $\alpha'>0$ a.e. in
$[-1,1]$, and $0<p<\infty$. There exists a constant $C>0$ such that if
$g$, $h$ are Lebesgue-measurable functions on $[-1,1]$, then
$$
   \Vert \alpha'(x)^{-1/2}(1-x^2)^{-1/4}g(x) \Vert _{L^p_*(\vert h \vert^pdx)}
   \leq C \liminf_{n \to \infty} \Vert P_ng \Vert _{L^p_*(\vert h \vert^pdx)}.
$$

\medskip

The following lemma will be useful to estimate some weighted $L^p_*$ norms:

\proclaim{Lemma 6}. Let $1 \leq p <\infty$, $r,s \in \R$, $a>0$. Then,
$$
   \chi_{(0,a)}(x) x^r \in L^p_*(x^sdx) \Longleftrightarrow
   pr+s+1 \geq 0, \ \ (r,s) \not= (0,-1).
$$
Moreover, in this case there is a constant $K$ depending on $r,s,p$ such that
$$
   \Vert \chi_{(0,a)}(x) x^r \Vert _{L^p_*(x^sdx)} = K a^{r+(s+1)/p}.
$$

\bigskip

{\noindent \bf \S 3. Proof of theorem 1.}

\medskip

\noindent The weak boundedness
$$
   \Vert uS_nf \Vert _{L^p_*(w)} \leq	 C \Vert uf \Vert _{L^p(w)}
$$
implies the following conditions (see [6], theorem 1, with the appropriate
changes):
$$
   u \in L^p_*(w),
$$
$$
   u^{-1} \in L^q(w),
$$
$$
   u(x) w(x)^{-1/2}(1-x^2)^{-1/4} \in L^p_*(w),
$$
$$
   u(x)^{-1} w(x)^{-1/2}(1-x^2)^{-1/4} \in L^q(w),
$$
where $1/p + 1/q = 1$. With the weight $u(x) = (1-x)^a(1+x)^b$ and having in
mind that $\alpha, \beta \geq -1/2$, this means
$$
   \eqalign{
   -{1 \over 4} \leq a+(\alpha+1)({1 \over p}-{1 \over 2}) < {1 \over 4}, \cr
   -{1 \over 4} \leq b+(\beta+1)({1 \over p}-{1 \over 2}) < {1 \over 4}.
   \cr}
$$
Therefore, we only need to show that the equality cannot occur in the left hand
side of these equations. Assume, for example,
$$
   -{1 \over 4} = a+(\alpha+1)({1 \over p}-{1 \over 2}). \eqno (7)
$$
Let us consider again Pollard's decomposition of the partial sums $S_nf$. We
will prove that there exists a constant $C$ such that
$$
   \Vert u W_{1,n}f \Vert _{L^p_*(w)} \leq	 C \Vert u f \Vert _{L^p(w)}
$$
and
$$
   \Vert u W_{3,n}f \Vert _{L^p(w)} \leq	 C \Vert u f \Vert _{L^p(w)}.
$$
This, together with the boundedness of $S_n$, implies the same property for
$W_{2,n}$ and will lead to a contradiction.

\medskip

a) {\it Boundedness of $W_{1,n}$}. From its definition, we have
$$
   \Vert u W_{1,n}f \Vert _{L^p_*(w)} \leq \Vert u P_{n+1} \Vert _{L^p_*(w)}
   \Vert u^{-1} P_{n+1} \Vert _{L^q(w)} \Vert u f \Vert _{L^p(w)}.
$$
So, we only need to prove
$$
   \Vert u P_n \Vert _{L^p_*(w)} \leq C \quad \forall n \in \N
$$
and
$$
   \Vert u^{-1} P_n \Vert _{L^q(w)} \leq C \quad \forall n \in \N,
$$
which follows from lemma 6, (5) and the dominate convergence theorem.

\medskip

b) {\it Boundedness of $W_{3,n}$}. Using again (5) and (6), it is enough
to obtain
$$
   \Vert Hg \Vert _{L^p(v)} \leq C \Vert g \Vert _{L^p(v)}
   \quad \forall g \in L^p(v),
$$
with
$$
   v(x) = (1-x)^{\alpha + a p + p(1 - 2\alpha)/4}
   (1+x)^{\beta + a p + p(1 - 2\beta)/4}.
$$
Now, we only need to prove that $v \in A_p$. This can be deduced from lemma 3.

\medskip

c) From $a)$, $b)$ and the hypothesis, we have a constant $C$ such that
for all $f \in L^p(u^pw)$ and every $n \in \N$
$$
   \Vert uW_{2,n}f \Vert _{L^p_*(w)} \leq	 C \Vert uf \Vert _{L^p(w)},
$$
that is,
$$
   \Vert u P_{n+1} Hg \Vert _{L^p_*(w)} \leq	 C
   \Vert u(x) (1-x^2)^{-1} Q_n(x)^{-1} w(x)^{-1} g \Vert _{L^p(w)}.
$$
Applying (4), we have
$$
   \Vert u P_{n+1} \chi_I \Vert _{L^p_*(w)} \left(
   \int_{-1}^1 { u(x)^{-q} (1-x^2)^q \vert Q_n(x) \vert^q w(x) \over
   (\vert I \vert + \vert x-x_I \vert)^q } dx
   \right)^{1/q} \leq C
$$
for every interval $I \subseteq [-1,1]$, with a constant $C>0$ independent of
$n$ and $I$; by lemma 5 with $I = [1-\varepsilon,1]$, it follows
$$
   \Vert x^{a-\alpha/2-1/4} \chi_{[0,\varepsilon]} \Vert _{L^p_*(x^\alpha)}
   \left( \int_0^1 { x^{-aq+q/4+\alpha(1-q/2)} \over
   (\varepsilon + \vert x-\varepsilon/2 \vert)^q } dx
   \right)^{1/q} \leq C. \eqno (8)
$$
Now, by lemma 6 and (7)
$$
   \Vert x^{a-\alpha/2-1/4} \chi_{[0,\varepsilon]} \Vert _{L^p_*(x^\alpha)}
   = K \eqno (9)
$$
and
$$
   \int_0^1 { x^{-aq+q/4+\alpha(1-q/2)} \over
   (\varepsilon + \vert x-\varepsilon/2 \vert)^q } dx =
   \int_0^1
   { x^{1/(p-1)} \over (\varepsilon + \vert x-\varepsilon/2 \vert)^q } dx \geq
   C \int_\varepsilon^1 x^{1/(p-1) - q} dx = C \vert \log \varepsilon \vert,
$$
which, together with (9), leads to a contradiction in (8). Therefore, (7)
cannot be true and the theorem is proved.

\bigskip

{\goodbreak\noindent \bf \S 4. Adding mass points.\nobreak}

\medskip

Let $d\mu$ be a positive measure on $\R$, $d\nu = d\mu + \sum_{i=1}^k M_i
\delta_{a_i}$, where $M_i > 0$, $\mu\{a_i\} = 0$. Let also $u$ be a weight
such that $0 < u(a_i) < \infty$ ($i=1,\dots,k$). We will denote by
$\{K_n(x,y)\}$ the kernels relative to $d\mu$ and by $\{L_n(x,y)\}$ the kernels
relative to $d\nu$. Then, the $n$-th partial sum of the Fourier series with
respect to $d\nu$ is given by
$$
   S_nf(x) = \int_{\R} L_n(x,y) f(y) d\nu(y).
$$
Let us take $1 < p < \infty$, $1/p + 1/q = 1$ and
$$
   T_nf(x) = \int_{\R} L_n(x,y) f(y) d\mu(y).
$$
Then

\proclaim{Theorem 7}. With the above notation, there exists a constant $C$ such
that
$$
   \Vert uS_nf \Vert _{L^p_*(d\nu)} \leq C \Vert uf \Vert _{L^p(d\nu)}
   \quad \forall n \geq 0, \ \forall f \in L^p(u^pd\nu) \eqno(10)
$$
if and only if there exists another constant $C$ such that:
\medskip
\item {a)} $\Vert uT_nf \Vert _{L^p_*(d\mu)} \leq C \Vert uf \Vert
_{L^p(d\mu)}\quad \forall n \geq 0, \ \forall f \in L^p(u^pd\mu)$;
\medskip
\item {b)} $u(a_i) \Vert u^{-1}L_n(x,a_i) \Vert _{L^q(d\mu)} \leq C
\quad \forall n \geq 0, \ (i=1,\dots,k)$;
\medskip
\item {c)} $\Vert uL_n(x,a_i) \Vert _{L^p_*(d\mu)} \leq C u(a_i)
\quad \forall n \geq 0, \ (i=1,\dots,k)$.
\medskip \noindent
The same holds replacing $L^p_*(d\nu)$ by $L^p(d\nu)$ and $L^p_*(d\mu)$
by $L^p(d\mu)$.

\proof From the definition, it follows
$$
   S_nf(x) = T_nf(x) + \sum_{i=1}^k M_iL_n(x,a_i)f(a_i). \eqno(11)
$$
Now, suppose $(10)$ holds. If $f \in L^p(u^pd\nu)$, let us define $g(x)=f(x)$
for $x \not= a_i$  $(i=1,\dots,k)$ and $g(a_i)=0$ $(i=1,\dots,k)$. Since
$\mu(\{a_i\})=0$, we have $S_ng=T_nf$ and
$$
   \Vert ug \Vert _{L^p(d\nu)} = \Vert uf \Vert _{L^p(d\mu)}.
$$
Therefore, (10) implies
$$
   \Vert uT_nf \Vert _{L^p_*(d\nu)} \leq C \Vert uf \Vert _{L^p(d\mu)}
   \quad \forall n \geq 0, \ \forall f \in L^p(u^pd\nu). \eqno(12)
$$
Taking now $f=\chi_{\{a_i\}}$ we obtain $S_nf(x) = M_iL_n(x,a_i)$ and $\Vert
uf \Vert _{L^p(d\nu)} = M_i^{1/p}u(a_i)$. Thus, (10) also implies
$$
   \Vert uL_n(x,a_i) \Vert _{L^p_*(d\nu)} \leq C u(a_i)
   \quad \forall n \geq 0, \ (i=1,\dots,k). \eqno(13)
$$
Actually, since $\Vert uf \Vert _{L^p(d\mu)}$, $u(a_i) \vert f(a_i) \vert$
$\leq$ $\Vert uf \Vert _{L^p(d\nu)}$ it is inmediate from $(11)$ that $(12)$ and
$(13)$ imply $(10)$. So, we only need to show that $(12)$ is equivalent to $a)$ 
and
$b)$ and that $(13)$ is the same as $c)$.

It is easy to see that
$$
   \Vert uL_n(x,a_i) \Vert ^p_{L^p_*(d\nu)}
   \leq \Vert uL_n(x,a_i) \Vert ^p_{L^p_*(d\mu)}
   + \sum_{j=1}^k M_ju(a_j)^p\vert L_n(a_j,a_i)\vert^p.
$$
Now, by Schwarz inequality we have
$$
   \vert L_n(a_j,a_i)\vert \leq L_n(a_j,a_j)^{1/2} L_n(a_i,a_i)^{1/2}
$$
and $\{L_n(a_i,a_i)\}_{n\geq 0}$ is a bounded sequence, since $\mu(\{a_i\})>0$.
Therefore
$$
   \Vert uL_n(x,a_i) \Vert ^p_{L^p_*(d\mu)}
   \leq \Vert uL_n(x,a_i) \Vert ^p_{L^p_*(d\nu)}
   \leq \Vert uL_n(x,a_i) \Vert ^p_{L^p_*(d\mu)} + C
$$
and $(13)$ is actually equivalent to $c)$.

Let us examine now condition $(12)$. It is easy to see that
$$
   \Vert uT_nf \Vert ^p_{L^p_*(d\nu)}
   \leq \Vert uT_nf \Vert ^p_{L^p_*(d\mu)}
   + \sum_{i=1}^k M_iu(a_i)^p\vert T_nf(a_i)\vert^p,
$$
$$
   \Vert uT_nf \Vert ^p_{L^p_*(d\mu)}
   \leq \Vert uT_nf \Vert ^p_{L^p_*(d\nu)}
$$
and
$$
   M_iu(a_i)^p\vert T_nf(a_i)\vert^p \leq \Vert uT_nf \Vert ^p_{L^p_*(d\nu)}.
$$
Thus, $(12)$ holds if and only if condition $a)$ holds together with
$$
   u(a_i) \vert T_nf(a_i)\vert \leq C \Vert uf \Vert _{L^p(d\mu)}
   \quad \forall n \geq 0 \ \forall f \in L^p(u^pd\mu) \ (i=1,\dots,k).
$$
Taking into acount that
$$
   T_nf(a_i) = \int_{\R} L_n(a_i,x)f(x)d\mu(x),
$$
this last inequality is simply $b)$.
\medskip\noindent
The proof can be rewritten with $L^p$ norms instead of $L^p_*$ norms.

\medskip

The operators $T_n$ can be handled in a similar way to expansions with respect
to $d\mu$. As to like in parts $b)$ and $c)$, let us introduce the following
notation:

\medskip

$d\mu^c(x) = (x-c)^2d\mu(x)$;

\medskip

$\{P_n^c\}$ is the sequence of orthonormal polynomials relative to $d\mu^c$;

\medskip

$P_n^c(x) = k_n^c x^n + \dots, \qquad k_n^c > 0$;

\medskip

$\{K_n^c(x,y)\}$ is the sequence of kernels relative to $d\mu^c$.

\medskip

Then (see [7])

\proclaim {Proposition 8}. Let $d\mu$ be a positive measure on $\R$, $c \in \R$,
$M>0$. Let $\{\widetilde P_n\}_{n\geq 0}$ be the polynomials orthonormal with
respect to $d\mu + M \delta_c$. Then, for each $n \in \N$ there exist two
constants $A_n, B_n \in (0,1)$ such that
$$
   \widetilde P_n(x) = A_n P_n(x) + B_n (x-c) P_{n-1}^c(x).
$$
Furthermore, if ${\rm supp\ }d\mu = [-1,1]$, $\mu ' > 0$ a.e. and $c \in
[-1,1]$, then
$$
   \lim_{n \to \infty} A_n K_{n-1}(c,c) = {1 \over \lambda(c)+M}
$$
and
$$
   \lim_{n \to \infty} B_n = {M \over \lambda(c)+M},
$$
where
$$
   \lambda(c) = \lim_{n \to \infty} {1 \over K_n(c,c)}.
$$

\medskip

We can also find some relations which involve the kernels.

\proclaim {Proposition 9}. Let $d\mu$ be a positive measure on $\R$, $c \in \R$
and $M>0$. Let $\{\widetilde K_n\}_{n\geq 0}$ be the kernels relative to $d\mu +
M\delta_c$. Then $\forall n \in \N$
$$
   \widetilde K_n(x,y) = {1 \over 1 + M K_n(c,c)} K_n(x,y)
   + {M K_n(c,c) \over 1 + M K_n(c,c)} (x-c)(y-c) K_{n-1}^c(x,y).
$$

\medskip

Propositions 8 and 9 lead to bounds for $\widetilde P_n$ and $\widetilde K_n$,
provided bounds for $P_n$, $P_n^c$, $K_n$, $K_n^c$ are known. These bounds,
together with theorem 7, can be used to prove theorem 2 (see [8]).

\bigskip

\centerline {\bf REFERENCES}

\bigskip

\frenchspacing

\item {{\bf [1]}} R. Askey and S. Wainger, {\it Mean convergence of expansions
in Laguerre and Hermite series}, Amer. J. Math. {\bf 87} (1965), 695--708.

\item {{\bf [2]}} {\rm F. V. Atkinson and W. N. Everitt, {\it Orthogonal
polynomials which satisfy second order differential equations}, in \lq \lq E.
B. Christoffel, the influence of his work on mathematics and the physical
sciences" (P. L. Butzer and F. Feh\'er, eds.), Birkh\"auser, 1981, pp.
173--181.}

\item {{\bf [3]}} V. M. Badkov, {\it Convergence in the mean and almost
everywhere of Fourier series in polynomials orthogonal on an interval},  Math.
USSR Sb. {\bf 24} (1974), 223--256.

\item {{\bf [4]}} T. S. Chihara, {\it Orthogonal polynomials and measures with
end point masses}, Rocky Mountain J. Math. {\bf 15} (1985), 705--719.

\item {{\bf [5]}} J. Garc\'\i a Cuerva and J. L. Rubio de Francia, {\it \lq \lq
Weighted norm inequalities and related topics"}, North-Holland, Amsterdam, 1985.

\item {{\bf [6]}} J. J. Guadalupe, M. P\'erez and J. L. Varona, {\it Weak
behaviour of Fourier-Jacobi series}, J. Approx. Theory {\bf 61} (1990),
222--238.

\item {{\bf [7]}} {\rm J. J. Guadalupe, M. P\'erez, F. J. Ruiz and J. L. Varona,
{\it Asymptotic behaviour of orthogonal polynomials relative to measures with
mass points} (preprint).}

\item {{\bf [8]}} {\rm J. J. Guadalupe, M. P\'erez, F. J. Ruiz and J. L. Varona,
{\it On mean and weak convergence of Fourier series relative to measures with
mass points} (preprint).}

\item {{\bf [9]}} R. Hunt, B. Muckenhoupt and R. Wheeden, {\it Weighted norm
inequalities for the conjugate function and Hilbert transform}, Trans. Amer.
Math. Soc. {\bf 176} (1973), 227--251.

\item {{\bf [10]}} T. Koornwinder, {\it Orthogonal polynomials with weight
function $(1-x)^{\alpha}(1+x)^{\beta}+M\delta(x+1)+N\delta(x-1)$}, Canad. Math.
Bull. {\bf 27} (1984), 205--214.

\item {{\bf [11]}} A. M. Krall, {\it Orthogonal polynomials satisfying fourth
order differential equations}, Proc. Roy. Soc. Edinburgh {\bf 87} (1981),
271--288.

\item {{\bf [12]}} L. L. Littlejohn, {\it On the classification of differential
equations having orthogonal polynomial solutions}, Ann. Mat. Pura Appl. {\bf
138} (1984), 35--53.

\item {{\bf [13]}} A. M\'at\'e, P. Nevai and V. Totik, {\it Asymptotics for the
ratio of the leading coefficients of orthogonal polynomials on the unit
circle}, Constr. Approx. {\bf 1} (1985), 63--69.

\item {{\bf [14]}} A. M\'at\'e, P. Nevai and V. Totik, {\it Necessary
conditions for weighted mean convergence of Fourier series in orthogonal
polynomials}, J. Approx. Theory {\bf 46} (1986), 314--322.

\item {{\bf [15]}} B. Muckenhoupt, {\it Mean convergence of Jacobi series},
Proc. Amer. Math. Soc. {\bf 23} (1969), 306--310.

\item {{\bf [16]}} B. Muckenhoupt, {\it Mean convergence of Hermite and
Laguerre series. I}, Trans. Amer. Math. Soc. {\bf 147} (1970), 419--431.

\item {{\bf [17]}} B. Muckenhoupt, {\it Mean convergence of Hermite and
Laguerre series. II}, Trans. Amer. Math. Soc. {\bf 147} (1970), 433--460.

\item {{\bf [18]}} B. Muckenhoupt and R. L. Wheeden, {\it Two weight function
norm inequalities for the Hardy-Littlewood maximal function and the Hilbert
transform}, Studia Math. {\bf 55} (1976), 279--294.

\item {{\bf [19]}} C. J. Neugebauer, {\it Inserting $A_p$ weights}, Proc. Amer.
Math. Soc. {\bf 87} (1983), 644--648.

\item {\bf [20]} {\rm P. Nevai, {\it \lq \lq Orthogonal Polynomials"}, Memoirs
of the Amer. Math. Soc., vol. 18, n. 213, Providence, RI, U.S.A., 1979.}

\item {{\bf [21]}} J. Newman and W. Rudin, {\it Mean convergence of orthogonal
series}, Proc. Amer. Math. Soc. {\bf 3} (1952), 219--222.

\item {{\bf [22]}} H. Pollard, {\it The mean convergence of orthogonal series.
II}, Trans. Amer. Math. Soc. {\bf 63} (1948), 355--367.

\item {{\bf [23]}} E. A. Rahmanov, {\it On the asymptotics of the ratio of
orthogonal polynomials. II}, Math. USSR. Sb. {\bf 46} (1983), 105--117.

\item {{\bf [24]}} G. M. Wing, {\it The mean convergence of orthogonal series},
Amer. J. Math. {\bf 72} (1950), 792--808.

\bye
